\section{Software Architecture Layers}
\label{sec:software-design}

Twelve layers, listed bottom-up, shared verbatim by both Option A and Option B. The two options differ only in \emph{how} layers 1--4 get populated --- Option A's Firmware Compiler emits a Machine IR entry directly at build time (Section~\ref{sec:e2b-family}); Option B's Correlation Engine and Graph-to-IR compiler emit one once a mapping clears a confidence threshold (Section~\ref{sec:probe-discovery}, \ref{sec:knowledge-graph}). Everything from layer 5 upward is identical code regardless of which path populated the entry below it, which is the mechanism behind the convergence argument in Section~\ref{sec:convergence}.

\begin{enumerate}
\item \textbf{Runtime} --- the process(es) executing on the ACC (and, in a thin form, on each E2B node) that host drivers, scheduling, and the network stack. On the ACC: a Linux-hosted supervisor process plus a set of containerized "skills"; on RT domain and E2B nodes: a small deterministic executive (Zephyr-class or bare-metal) running generated driver code.
\item \textbf{Hardware Abstraction Layer (HAL)} --- per-interface-type drivers (PWM, ADC, SPI, CAN-FD, LIN, SENT, Ethernet) with a uniform internal API (\code{hal\_pwm\_set\_duty()}, \code{hal\_can\_send\_frame()}, etc.) so that the Verified Primitive Library above never touches a register directly.
\item \textbf{Verified Primitive Library} --- a curated set of driver \emph{instances} for specific hardware (a specific motor driver IC, a specific Hall sensor part, a specific CAN transceiver's init sequence) that have been through the verification pipeline in Section~\ref{sec:firmware-generation}. This is the layer that actually grows as a moat (Section~\ref{sec:moat}): each new peripheral supported adds one verified entry here, reusable on every future machine.
\item \textbf{Machine IR} (Section~\ref{sec:machine-ir}) --- the structured, versioned model of a specific machine instance: which device instances exist, which network path/resource binding each one uses, and (for legacy devices) which mapping confidence applies.
\item \textbf{Hardware Knowledge Graph} --- the cross-machine, cross-program database of \emph{general} facts: "this MCU part has N timer channels usable for PWM," "this transceiver part requires this init sequence," "this CAN ID pattern was previously seen meaning turn-signal on a different vehicle from the same OEM platform." The Machine IR for one vehicle references entries in this graph; the graph itself is shared across every vehicle Alfred has ever touched.
\item \textbf{Hardware Mapping / Device Abstraction} --- the layer that resolves a semantic capability call to a specific Machine-IR device instance and its bound primitive(s). This is literally the "Alfred hardware mapping" box in the architecture diagram in the background brief.
\item \textbf{Alfred Semantic API} --- the stable, typed verb surface (\code{door.open()}, \code{actuator.set\_position()}) that application code is written against.
\item \textbf{Behavior / Application Layer} --- the actual vehicle behaviors and business logic (e.g., "when door opens, if vehicle speed $>$ 0, ignore"), composed from Semantic API calls. This is where OEM/customer-specific logic lives, deliberately kept separate from everything below it.
\item \textbf{Firmware Compiler} (Section~\ref{sec:firmware-generation}) --- the deterministic toolchain that turns a Machine IR + placement decisions into buildable, flashable firmware images for the ACC RT domain and every E2B node.
\item \textbf{Target Packs} --- per-MCU/per-PHY "backend" definitions for the Firmware Compiler (pin maps, timer/DMA resource tables, linker scripts, board support package) --- adding a new MCU family to Alfred means writing one target pack, not touching the compiler.
\item \textbf{Deployment Manager} (Section~\ref{sec:deployment-reprogramming}) --- orchestrates SIL/HIL verification gates, staged rollout, and the actual programming of ACC and E2B targets, with rollback.
\item \textbf{Diagnostics \& Certification Evidence System} --- structured logging of every decision the pipeline made (resource allocation choices, test results, static analysis reports) in a form that is directly usable as functional-safety work-product evidence (ISO~26262 Part 6/8 traceability), not just human-readable debug output.
\end{enumerate}

\begin{diagram}
   Behavior / Application Layer  <-- OEM/customer logic; unchanged across A/B
              |
      Alfred Semantic API        <-- stable verbs
              |
    Hardware Mapping / Device Abstraction
              |
        +-----+------------------------+
        |                               |
   Machine IR (per-vehicle)     Hardware Knowledge Graph (cross-vehicle)
        |                               |
        +-------------+-----------------+
                      |
         Verified Primitive Library  <-- grows every program, reused forever
                      |
                     HAL
                      |
                   Runtime
\end{diagram}

Option B's population path additionally requires a small set of software surfaces Option A never needs: a Capture Service (per probe node), a Protocol Identification Engine (deterministic Level 1/2 decoders), a Correlation Engine (the Level 3 statistical machinery in Section~\ref{sec:probe-discovery}), an Engineering File Importer (DBC/ARXML/LDF/A2L/ODX), and the Vehicle Knowledge Graph builder (Section~\ref{sec:knowledge-graph}) that compiles a confidence-cleared graph subtree into a Machine IR entry. These are the only genuinely new software components Option B requires beyond what Option A already needs --- everything from layer 5 upward is shared, unmodified code.
